\documentclass[11pt,a4paper]{article}
\usepackage[utf8]{inputenc}
\usepackage[spanish]{babel}
\usepackage[top=2cm,bottom=2cm,left=2cm,right=2cm]{geometry}
\usepackage{booktabs}
\usepackage{xcolor}
\usepackage{colortbl}
\usepackage{array}
\usepackage{titlesec}
\usepackage{mdframed}
\usepackage{enumitem}
\usepackage{listings}
\usepackage{hyperref}

\definecolor{techblue}{RGB}{0,86,179}
\definecolor{lightgray}{RGB}{245,245,245}
\definecolor{darkgray}{RGB}{80,80,80}
\definecolor{codegreen}{RGB}{40,120,40}

% Configuración de listings para código
\lstset{
  backgroundcolor=\color{lightgray},
  basicstyle=\small\ttfamily,
  keywordstyle=\color{techblue}\bfseries,
  commentstyle=\color{darkgray}\itshape,
  stringstyle=\color{codegreen},
  numbers=none,
  breaklines=true,
  frame=single,
  rulecolor=\color{lightgray!80!black},
  showstringspaces=false,
  keepspaces=true
}

\titleformat{\section}{\bfseries\color{techblue}\large}{}{0em}{}[{\vspace{-3pt}\color{techblue}\hrule height 0.5pt}\vspace{3pt}]
\titlespacing{\section}{0pt}{12pt}{6pt}

\setlength{\parindent}{0pt}
\setlength{\parskip}{6pt}

\begin{document}

\begin{center}
  {\Large\bfseries\color{techblue}Guía de Configuración: Logitech M720 Triathlon}\\[4pt]
  {\small Controlador LogiOps (logid) en Ubuntu 24.04 LTS (Wayland)}
\end{center}

\vspace{5pt}

\begin{mdframed}[backgroundcolor=techblue!5,linecolor=techblue,linewidth=1pt,
    innertopmargin=6pt,innerbottommargin=6pt,innerleftmargin=8pt,innerrightmargin=8pt]
  {\small\textbf{Nota del sistema:} Esta guía detalla la instalación, configuración y depuración del daemon de Logitech \texttt{logid} para habilitar funciones especiales en el mouse Logitech M720 Triathlon, tales como gestos con el pulgar y asignaciones de teclado personalizadas.}
\end{mdframed}

\section{1. Instalación de Dependencias}
Para poder compilar e instalar el controlador en Ubuntu 24.04, es necesario contar con las siguientes dependencias de desarrollo. Instálalas con el siguiente comando:

\begin{lstlisting}[language=bash]
sudo apt update && sudo apt install -y build-essential cmake pkg-config libevdev-dev libudev-dev libconfig++-dev libglib2.0-dev
\end{lstlisting}

\section{2. Compilación e Instalación}
Una vez obtenidas las dependencias, realiza la compilación del código desde el directorio del repositorio clonado:

\begin{lstlisting}[language=bash]
cd ~/src/trabajo_de_titulo/rebaja_arancel/logiops
mkdir -p build && cd build
cmake -DCMAKE_BUILD_TYPE=Release ..
make
sudo make install
\end{lstlisting}

Este proceso generará el ejecutable binario e instalará el archivo del servicio de systemd para controlar el daemon.

\section{3. Archivo de Configuración}
El daemon busca su archivo de configuración en la ruta global \texttt{/etc/logid.cfg}. Crea y edita este archivo con privilegios de root:

\begin{lstlisting}[language=bash]
sudo nano /etc/logid.cfg
\end{lstlisting}

A continuación se muestra un ejemplo optimizado para el mouse \textbf{Logitech M720 Triathlon}, el cual mapea el botón del pulgar para gestos del sistema y control multimedia:

\begin{lstlisting}
devices: (
{
    name: "Logitech M720 Triathlon";
    dpi: 1000; // Ajuste de sensibilidad predeterminada

    buttons: (
        // Boton oculto bajo el pulgar (Gesture button)
        {
            cid: 0xc3;
            action =
            {
                type: "Gestures";
                gestures: (
                    {
                        direction: "Up";
                        mode: "OnRelease";
                        action = {
                            type: "Keypress";
                            keys: ["KEY_PLAYPAUSE"]; // Mover arriba: Play/Pausa
                        };
                    },
                    {
                        direction: "Down";
                        mode: "OnRelease";
                        action = {
                            type: "Keypress";
                            keys: ["KEY_MUTE"]; // Mover abajo: Silenciar
                        };
                    },
                    {
                        direction: "Left";
                        mode: "OnRelease";
                        action = {
                            type: "Keypress";
                            keys: ["KEY_LEFTMETA", "KEY_LEFT"]; // Mover izq.: Escritorio izquierdo
                        };
                    },
                    {
                        direction: "Right";
                        mode: "OnRelease";
                        action = {
                            type: "Keypress";
                            keys: ["KEY_LEFTMETA", "KEY_RIGHT"]; // Mover der.: Escritorio derecho
                        };
                    },
                    {
                        direction: "None";
                        mode: "OnRelease";
                        action = {
                            type: "Keypress";
                            keys: ["KEY_LEFTMETA"]; // Clic simple: Vista de Actividades
                        };
                    }
                );
            };
        }
    );
}
);
\end{lstlisting}

\section{4. Habilitar y Controlar el Servicio}
Para aplicar los cambios e iniciar la ejecución automática del daemon en segundo plano, controla el servicio con los comandos de \texttt{systemctl}:

\begin{lstlisting}[language=bash]
# Habilitar e iniciar por primera vez
sudo systemctl enable --now logid

# Detener el servicio (por ejemplo, para pruebas)
sudo systemctl stop logid

# Reiniciar el servicio (tras hacer cambios en la configuracion)
sudo systemctl restart logid
\end{lstlisting}

\section{5. Detección de Botones y Depuración (CIDs)}
Si deseas mapear otros botones del mouse, debes identificar su identificador numérico o \texttt{cid}. Para hacerlo en tiempo real:

\begin{enumerate}
    \item Detén el servicio en ejecución: \texttt{sudo systemctl stop logid}.
    \item Ejecuta el binario manualmente en modo detallado: \texttt{sudo logid -v}.
    \item Presiona los botones del mouse. La terminal mostrará los eventos con su correspondiente código CID (por ejemplo, \texttt{cid: 0x53} para retroceder).
    \item Presiona \texttt{Ctrl + C} en la terminal para finalizar las pruebas y reinicia el servicio con \texttt{sudo systemctl start logid}.
\end{enumerate}

\end{document}
